\clearpage\thispagestyle{empty}
\pdfbookmark[0]{Quellen zum Kometenbericht vom Juni 1824}{cometjun1824-sources}
\begin{center}\textbf{Quellen zum Kometenbericht vom Juni 1824}\end{center}
\noindent Diese KI-Übertragung ins heutige Deutsch folgt dem Bericht \textit{Comet} in \href{https://gdz.sub.uni-goettingen.de/id/PPN236007467}{den \textit{Werken}, Band~VI} (Göttingen, 1874), S.~450. Gauß’ Bericht erschien zuerst unter der Überschrift \textit{Beobachtungen des Cometen von 1824 in Göttingen} in den \textit{Astronomischen Nachrichten}, Band~III, Nr.~59 (Juni 1824), \href{https://api.digitale-sammlungen.de/iiif/image/v2/bsb10049428_00102/full/1800,/0/default.jpg}{Sp.~179--180} und \href{https://api.digitale-sammlungen.de/iiif/image/v2/bsb10049428_00103/full/1800,/0/default.jpg}{Sp.~181--182}. Die Links führen zum Exemplar der Bayerischen Staatsbibliothek. H.~C. Schumacher gab den Band 1825 in Altona heraus.
\par\smallskip\noindent Der Erstdruck gibt als mittlere Zeit für den 3. Februar 11 55 7,1 an und endet mit Gauß’ Unterschrift. Die Werkausgabe gibt 11 35 7.1 an und endet mit der letzten Beobachtung. Die Übertragung folgt der Werkausgabe und behält die Leerstelle und die vier Punkte am 20. Februar sowie den abgesetzten Punkt vor der Rektaszension am 27. Februar bei. KI-Übertragung und KI-Satz.
\clearpage
